% =============================================================================
% meetAI Graduation Project Thesis
% Preamble: encoding, fonts, page layout, packages
% =============================================================================

\documentclass[12pt,a4paper,oneside]{book}

% --- Encoding and fonts ---
\usepackage[utf8]{inputenc}
\usepackage[T1]{fontenc}
\usepackage{mathpazo}             % Palatino text + matching math
\linespread{1.05}                 % Palatino benefits from a bit more leading
\renewcommand{\ttdefault}{lmtt}  % Latin Modern Mono for code

% --- Page layout ---
\usepackage[a4paper, margin=1in, bindingoffset=0.5cm]{geometry}
\usepackage[final,protrusion=true,expansion=true]{microtype}
\usepackage{setspace}
\onehalfspacing
\usepackage{emptypage}

% --- Math ---
\usepackage{mathtools}
\usepackage{amssymb}

% --- Graphics and figures ---
\usepackage{graphicx}
\usepackage[font=small,labelfont=bf,format=hang]{caption}
\usepackage{subcaption}
\usepackage{float}

% --- TikZ for in-document diagrams ---
\usepackage{tikz}
\usetikzlibrary{shapes.geometric, arrows.meta, positioning, fit, calc, backgrounds}

% --- Tables ---
\usepackage{booktabs}
\usepackage{array}
\usepackage{multirow}
\usepackage{tabularx}
\usepackage{longtable}

% --- Lists ---
\usepackage{enumitem}

% --- Code listings (Python / TS / SQL) ---
\usepackage{listings}
\lstset{
  basicstyle=\ttfamily\footnotesize,
  keywordstyle=\color{blue!70!black}\bfseries,
  commentstyle=\color{gray}\itshape,
  stringstyle=\color{red!60!black},
  showstringspaces=false,
  breaklines=true,
  frame=single,
  framerule=0.4pt,
  numbers=left,
  numberstyle=\tiny\color{gray},
  tabsize=2,
}

% --- Colors ---
\usepackage{xcolor}
\definecolor{chaptercolor}{RGB}{0, 70, 140}
\definecolor{linkblue}{RGB}{0,51,153}
\definecolor{tbdcolor}{RGB}{220, 80, 30}

% --- Header/Footer ---
\usepackage{fancyhdr}
\pagestyle{fancy}
\fancyhf{}
\fancyhead[L]{\small\nouppercase{\leftmark}}
\fancyhead[R]{\small\thepage}
\renewcommand{\headrulewidth}{0.4pt}
\fancypagestyle{plain}{
    \fancyhf{}
    \fancyfoot[C]{\small\thepage}
    \renewcommand{\headrulewidth}{0pt}
}

% --- Chapter title styling ---
\usepackage{titlesec}
\titleformat{\chapter}[display]
  {\normalfont\huge\bfseries\color{chaptercolor}}
  {\chaptertitlename\ \thechapter}{20pt}{\Huge}
\titlespacing*{\chapter}{0pt}{-20pt}{40pt}

\titleformat{\section}
  {\normalfont\Large\bfseries\color{chaptercolor}}
  {\thesection}{1em}{}
\titleformat{\subsection}
  {\normalfont\large\bfseries\color{chaptercolor!85!black}}
  {\thesubsection}{1em}{}

% --- Algorithms ---
% Note: the algorithm/algpseudocode packages are not in the base TeX Live
% install on this machine. We provide a self-contained `algorithm` float and
% a `listofalgorithms` is generated by the front matter (manual entry).
\newcounter{algorithm}
\renewcommand{\thealgorithm}{\thechapter.\arabic{algorithm}}
\newenvironment{algorithm}[2][htb]
  {\refstepcounter{algorithm}\par\noindent\textbf{Algorithm \thealgorithm: #2}\par\addcontentsline{loa}{section}{Algorithm \thechapter.\thealgorithm: #2}\par\noindent\rule{\linewidth}{0.4pt}\par\smallskip}
  {\par\smallskip\noindent\rule{\linewidth}{0.4pt}}

% --- Theorems ---
\usepackage{amsthm}
\theoremstyle{plain}
\newtheorem{theorem}{Theorem}[chapter]
\newtheorem{lemma}[theorem]{Lemma}
\theoremstyle{definition}
\newtheorem{definition}[theorem]{Definition}
\newtheorem{example}[theorem]{Example}
\theoremstyle{remark}
\newtheorem{remark}[theorem]{Remark}

% --- Units and numbers ---
% siunitx is in texlive-science (not installed on this machine). Raw SI
% units are written inline in the text instead (e.g., \SI{95}{\percent} is
% written as 95\,\%).

% --- Quotations ---
\usepackage[autostyle]{csquotes}

% --- Citations (natbib, IEEE numerical style) ---
\usepackage[numbers, sort&compress]{natbib}
\bibliographystyle{ieeetr}

% --- Appendices ---
\usepackage{appendix}

% --- Hyperlinks ---
\usepackage{bookmark}
\usepackage{hyperref}
\hypersetup{
    colorlinks=true,
    linkcolor=chaptercolor,
    citecolor=linkblue,
    urlcolor=linkblue,
    pdfauthor={meetAI Graduation Project Team},
    pdftitle={meetAI: Realtime Code-Switched Meeting Assistant with Speaker-Aware Retrieval-Augmented QA},
    pdfsubject={Graduation Project Thesis},
    pdfkeywords={ASR, speaker diarization, RAG, code-switched, Egyptian Arabic, meetAI},
    bookmarks=true,
    bookmarksnumbered=true,
    bookmarksopen=true,
}

% --- Smart cross-references ---
\usepackage{cleveref}

% --- TBD placeholder macro (for tables/metrics) ---
\newcommand{\tbd}[1][TBD]{\textcolor{tbdcolor}{\textbf{[#1]}}}
\newcommand{\todo}[1]{\textcolor{tbdcolor}{\textbf{[TODO: #1]}}}

% --- Custom commands ---
\newcommand{\R}{\mathbb{R}}
\newcommand{\N}{\mathbb{N}}
\DeclareMathOperator*{\argmax}{arg\,max}
\DeclareMathOperator*{\argmin}{arg\,min}
\DeclarePairedDelimiter{\abs}{\lvert}{\rvert}
\DeclarePairedDelimiter{\norm}{\lVert}{\rVert}

% --- Project-specific shorthand ---
\newcommand{\meetai}{meetAI}
\newcommand{\wer}{\mathrm{WER}}
\newcommand{\cer}{\mathrm{CER}}
\newcommand{\der}{\mathrm{DER}}
\newcommand{\mrr}{\mathrm{MRR}}
\newcommand{\ndcg}{\mathrm{nDCG}}

\begin{document}

% =============================================================================
% FRONT MATTER
% =============================================================================
\frontmatter

% --- Title page ---
\begin{titlepage}
\begin{center}
\vspace*{1.5cm}

{\Large\textbf{University Name}}\\[0.3cm]
{\large Faculty of Computers and Artificial Intelligence}\\[0.2cm]
{\large Department of Computer Science}\\[2.5cm]

{\Huge\bfseries\color{chaptercolor}
\meetai{}: A Realtime Meeting Assistant\\
for Code-Switched Egyptian Arabic\\[0.4cm]
with Speaker-Aware Retrieval-Augmented QA}\\[2cm]

{\Large\textit{A graduation project thesis submitted in partial fulfillment\\
of the requirements for the degree of}}\\[0.8cm]

{\Large\bfseries Bachelor of Science in Computer Science}\\[2cm]

{\large\textbf{Author:}}\\[0.3cm]
{\large Author One \quad Author Two \quad Author Three \quad Author Four}\\[1cm]

{\large\textbf{Supervisor:}}\\[0.3cm]
{\large Supervisor Name}\\[2cm]

{\large Academic Year 2025/2026}

\vfill
\end{center}
\end{titlepage}

% --- Declaration of Originality ---
\chapter*{Declaration of Originality}
\addcontentsline{toc}{chapter}{Declaration of Originality}

We, the undersigned, declare that this graduation project thesis represents our
original work, and that all sources of information have been duly acknowledged.
Where the work of others has been consulted (including, but not limited to,
open-source codebases, research papers, and software libraries), this is
clearly cited in the text and acknowledged in the bibliography. This thesis
has not been submitted, in whole or in part, for any other degree or
qualification.

\vspace{2cm}
\noindent\rule{6cm}{0.4pt}\\
Author One, \today

\noindent\rule{6cm}{0.4pt}\\
Author Two, \today

\noindent\rule{6cm}{0.4pt}\\
Author Three, \today

\noindent\rule{6cm}{0.4pt}\\
Author Four, \today

% --- Abstract ---
\chapter*{Abstract}
\addcontentsline{toc}{chapter}{Abstract}

\meetai{} is a real-time meeting assistant that captures live speech, performs
speaker-attributed transcription, and continuously updates a retrieval-augmented
knowledge base that users can query during and after the meeting. The system's
primary technical contribution is a fine-tuned automatic speech recognition
(ASR) model for code-switched Egyptian Arabic--English speech, paired with a
local-agreement word-by-word streaming commit policy that yields low-latency
yet stable partial transcripts. Speaker diarization produces per-segment
speaker labels, and a temporal-aware question-answering (QA) agent can resolve
natural-language time references (e.g., ``What did Speaker~1 say in the last
ten minutes?'') into structured queries over the meeting knowledge base. The
work draws on recent advances in large-scale ASR~\citep{radford2023whisper},
self-attention architectures~\citep{vaswani2017attention}, low-rank
parameter-efficient fine-tuning~\citep{hu2021lora}, streaming
diarization~\citep{guo2024sortformer}, streaming ASR commit
policies~\citep{liu2020streaming}, and retrieval-augmented
generation~\citep{lewis2020rag, johnson2019faiss}.

The backend is a polyglot microservices stack: a Node.js gateway, an
authentication service, a meeting-lifecycle service, and a Python AI engine
that exposes WebSocket-based streaming ASR, an embedder service, and a
server-sent-event (SSE) QA service. PostgreSQL with the pgvector extension
serves as the primary store, while Redis handles ephemeral state (token
indices, Redis Streams for the audio pipeline) and Anthropic's Gemini model
is used for topic summarization and answer generation. The system is
containerized with Docker and orchestrated via Docker Compose.

This thesis documents the problem definition, literature review, requirements
engineering, system architecture, machine-learning methodology (including
data collection, low-rank adaptation of Whisper, and the local-agreement
streaming algorithm), experimental design, and the resulting evaluation
metrics. The contributions of this work are: (i)~a curated code-switched
Egyptian Arabic speech corpus, (ii)~a LoRA-fine-tuned ASR model with
local-agreement streaming, (iii)~a realtime retrieval-augmented knowledge
base, and (iv)~a temporal-speaker-aware QA agent for meeting transcripts.

% --- Acknowledgments ---
\chapter*{Acknowledgments}
\addcontentsline{toc}{chapter}{Acknowledgments}

We would like to express our gratitude to our project supervisor for the
guidance, patience, and feedback throughout this work. We thank our families
and colleagues for their continuous support, and the open-source community
for the tools, models, and datasets that made this project possible.

% --- Dedication (optional) ---
\chapter*{Dedication}
\addcontentsline{toc}{chapter}{Dedication}

To our families, supervisors, and everyone who believed in the value of
realtime, language-aware, and accessible meeting technology.

% --- Table of contents and lists ---
\tableofcontents

\listoffigures
\listoftables

% The list of algorithms is collected automatically as `algorithm` floats
% are added throughout the chapters. A manual section is provided here to
% satisfy the structural framework, and will populate as the chapters are
% written.
\chapter*{List of Algorithms}
\addcontentsline{toc}{chapter}{List of Algorithms}

A specific index of all numbered algorithms used in this thesis. As
algorithms are introduced in subsequent chapters they will be
auto-collected by the \texttt{loa} (list of algorithms) channel. Until
then this list will be empty; once a chapter introduces its first
algorithm (e.g., \texttt{Algorithm 5.1: Local-Agreement Commit}), it
will appear here.

% Equations and abbreviations are inline throughout, but a dedicated list is
% generated here per the framework guide.
\chapter*{List of Equations}
\addcontentsline{toc}{chapter}{List of Equations}

A catalogue of mathematical models and loss functions used throughout this
thesis is auto-generated by the equation index. (Equations are numbered
per-chapter as \texttt{(\thechapter.\arabic{equation})}.)

\chapter*{List of Abbreviations and Acronyms}
\addcontentsline{toc}{chapter}{List of Abbreviations and Acronyms}

\begin{description}[itemsep=0.2em,leftmargin=3.5cm,style=nextline]
  \item[ANN] Approximate Nearest Neighbor
  \item[API] Application Programming Interface
  \item[ASR] Automatic Speech Recognition
  \item[CI/CD] Continuous Integration / Continuous Deployment
  \item[CPU] Central Processing Unit
  \item[DER] Diarization Error Rate
  \item[EM] Expectation Maximization
  \item[FAQ] Frequently Asked Questions
  \item[GPU] Graphics Processing Unit
  \item[JSONB] JSON Binary (PostgreSQL data type)
  \item[KB] Knowledge Base
  \item[LLM] Large Language Model
  \item[LoRA] Low-Rank Adaptation
  \item[MRR] Mean Reciprocal Rank
  \item[nDCG] normalized Discounted Cumulative Gain
  \item[NER] Named Entity Recognition
  \item[PCM] Pulse-Code Modulation
  \item[PEFT] Parameter-Efficient Fine-Tuning
  \item[QA] Question Answering
  \item[RAG] Retrieval-Augmented Generation
  \item[REST] Representational State Transfer
  \item[RMSE] Root Mean Square Error
  \item[RTF] Real-Time Factor
  \item[SDK] Software Development Kit
  \item[SSE] Server-Sent Events
  \item[SOTA] State of the Art
  \item[SQL] Structured Query Language
  \item[UI] User Interface
  \item[VAD] Voice Activity Detection
  \item[WER] Word Error Rate
  \item[CER] Character Error Rate
  \item[WS] WebSocket
\end{description}

% --- Glossary ---
\chapter*{Glossary}
\addcontentsline{toc}{chapter}{Glossary}

\begin{description}[itemsep=0.2em,leftmargin=3.5cm,style=nextline]
  \item[Code-switching] The alternation between two or more languages within a
  single conversation or even a single utterance. In this thesis we focus on
  Egyptian Arabic--English code-switching.
  \item[Commit policy] A streaming ASR decision rule for finalizing partial
  hypotheses. The \emph{local-agreement} policy emits a word when the model
  has produced it in at least $k$ consecutive partial decodes.
  \item[Diarization] The process of partitioning an audio stream into
  speaker-homogeneous segments and labeling them with a speaker identity.
  \item[Embedding] A dense real-valued vector representation of text (or audio)
  used for similarity search and retrieval.
  \item[Hallucination] The phenomenon in which a generative model produces
  fluent but factually incorrect content.
  \item[pgvector] A PostgreSQL extension that adds the \texttt{vector} data
  type, ANN indexes (ivfflat / HNSW), and cosine / L2 distance operators.
  \item[Speaker map] A user-managed mapping from system-assigned speaker labels
  (\texttt{spk\_0}, \texttt{spk\_1}, \ldots) to human-readable identities
  (e.g., ``Alice'').
  \item[Token-level streaming] The practice of producing incremental
  transcription output at the token (sub-word) granularity rather than
  waiting for an utterance to finish.
\end{description}

% --- List of publications ---
\chapter*{List of Publications}
\addcontentsline{toc}{chapter}{List of Publications}

To be added upon submission / acceptance of any related peer-reviewed
publications, posters, or conference presentations derived from this work.

% =============================================================================
% MAIN MATTER
% =============================================================================
\mainmatter

% --- Chapter files are included from chapters/ ---
% =============================================================================
% Chapter 1: Introduction
% =============================================================================

\chapter{Introduction}
\label{ch:intro}

Meetings are the backbone of how organizations communicate,
 yet the tools built to support them have not kept pace with how people actually speak.
In Egypt, business conversations naturally mix Egyptian Arabic and English within the same sentence.
This is not a communication problem; it is simply how professionals here talk.
The problem is that every mainstream meeting assistant on the market either struggles with Arabic entirely,
forces a fallback to Modern Standard Arabic that nobody uses in a boardroom,
or processes speech after the meeting is over. None of them can answer the question
عمر قال ايه لميار في أول الـ stand up meeting?" while the meeting is still running. This thesis builds a system that can.

\section{Problem Statement}
\label{sec:problem-statement}

We define the problem addressed in this thesis as follows. Given a live
multilingual meeting (with Egyptian Arabic--English code-switching as the
primary scenario), we must: (1)~capture audio in real time, (2)~produce a
low-latency word-by-word transcript that updates visibly to participants, (3)~attribute each spoken segment to a
distinct speaker, (4)~continuously index the emerging transcript into a
retrieval-augmented knowledge base, and (5)~answer natural-language questions
that refer to specific speakers and time windows, e.g., \textit{``هو أيمن قال ايه عن الـ budget اخر عشر دقايق؟''} The combined requirement of
low-latency streaming, code-switched accuracy, and temporal-speaker-aware
retrieval is the central technical problem.

\section{Motivation and Significance}
\label{sec:motivation}

Several real-world pain points motivate this work. Professional meeting
participants spend an estimated 30--50\% of their cognitive load on
note-taking and recall, which directly reduces their ability to engage
with the discussion. Post-meeting, the lack of a structured, speaker-
attributed transcript causes lost decisions and disputes about
responsibility. In multilingual regions such as the MENA area,
general-purpose ASR systems either fail on Egyptian Arabic phonology or
fall back to Modern Standard Arabic, neither of which matches natural
code-switched speech.

The significance of solving this problem in a single integrated system is
threefold. First, it produces measurable productivity gains for individual
users by removing the need for active note-taking. Second, it produces
measurable team-level gains by generating a structured, queryable
artifact from every meeting. Third, by targeting code-switched Egyptian
Arabic specifically, the work addresses an underserved language community
in commercial speech technology.

\section{Research Gaps, Questions, and Objectives}
\label{sec:rqs}

The literature (reviewed in \cref{ch:background}) contains strong solutions
to most individual sub-problems: large-vocabulary ASR via Whisper-family
models, streaming ASR with chunk-and-commit policies, speaker
diarization via Sortformer or similar architectures, and retrieval-
augmented generation (RAG) over vector stores. However, the integrated
problem stated in \cref{sec:problem-statement} contains a research gap:
code-switched Egyptian Arabic ASR with a low-latency stable commit policy
is not addressed by off-the-shelf systems, and temporal-speaker-aware
querying of a realtime RAG index is not natively supported by existing
meeting-assistant tools. The specific research questions addressed in this
thesis are:

\begin{enumerate}[label=\textbf{RQ\arabic*.}, itemsep=0.4em]
  \item \textbf{RQ1 (Streaming ASR for code-switched Egyptian Arabic).} To
  what extent can parameter-efficient fine-tuning (LoRA) of a Whisper-
  family model, combined with a local-agreement commit policy, produce
  realtime transcripts whose Word Error Rate (WER) on a held-out
  code-switched test set is below an operational threshold (TBD) while
  maintaining sub-second first-token latency?
  \item \textbf{RQ2 (Speaker attribution and temporal QA).} To what extent
  does combining speaker-diarization output with a temporal-aware QA agent
  allow natural-language questions that reference specific speakers and
  time windows to be answered with relevant evidence (Recall@$k$ TBD)?
  \item \textbf{RQ3 (Realtime knowledge-base freshness).} What is the end-
  to-end latency between a spoken utterance and its retrievability via the
  QA system, and what chunking / topic-windowing strategy minimizes this
  latency while preserving retrieval quality (MRR TBD)?
  \item \textbf{RQ4 (System robustness).} Does the integrated system
  remain stable (no transcript regressions, no speaker label swaps) over
  sessions of at least 60 minutes of continuous speech, and what is the
  observed failure mode distribution?
\end{enumerate}

\subsection{Objectives}

The objectives derived from the research questions are:

\begin{itemize}[itemsep=0.3em]
  \item \textbf{O1.} Curate a code-switched Egyptian Arabic--English speech
  corpus with validated transcripts and per-speaker labels.
  \item \textbf{O2.} Train and validate a LoRA-fine-tuned ASR model that
  outperforms the off-the-shelf baseline on the curated test set.
  \item \textbf{O3.} Implement a local-agreement streaming commit policy
  that produces stable word-by-word output with configurable latency.
  \item \textbf{O4.} Integrate a pretrained speaker-diarization model with
  the streaming ASR pipeline to produce speaker-attributed transcripts.
  \item \textbf{O5.} Implement a realtime RAG pipeline that ingests
  transcript chunks and supports temporal-speaker-aware retrieval.
  \item \textbf{O6.} Build and validate the full integrated system and
  report quantitative evaluation results.
\end{itemize}

\section{Proposed Solution Overview}
\label{sec:solution-overview}

The proposed solution, \meetai{}, is a realtime meeting assistant that
combines: (i)~a fine-tuned streaming ASR model for code-switched Egyptian
Arabic, (ii)~a local-agreement commit policy that finalizes words when they
appear in a configurable number of consecutive partial decodes, (iii)~a
pretrained Sortformer-based streaming speaker-diarization model, (iv)~a
realtime RAG knowledge base backed by PostgreSQL with the pgvector
extension, and (v)~a temporal-speaker-aware QA agent that parses natural-
language time and speaker references into structured retrieval queries. The
system is exposed through a thin React frontend and a polyglot
microservices backend (Node.js + Python).

\Cref{fig:high-level} shows the high-level system context. The full
component architecture is presented in \cref{ch:architecture}.

\begin{figure}[H]
\centering
\begin{tikzpicture}[
  node distance=1.2cm and 1.6cm,
  every node/.style={font=\small},
  box/.style={rectangle, draw, rounded corners, minimum width=3.2cm, minimum height=1cm, align=center, fill=blue!5},
  ext/.style={rectangle, draw, rounded corners, minimum width=3.2cm, minimum height=1cm, align=center, fill=orange!10},
  arrow/.style={-{Stealth[length=2mm]}, thick}
]
\node[ext] (user) {End user\\(browser)};
\node[box, right=of user] (fe) {React\\Frontend};
\node[box, right=of fe] (gw) {Node.js\\Gateway};
\node[box, right=of gw, yshift=0.8cm] (asr) {AI Engine\\(ASR + diarization)};
\node[box, right=of gw, yshift=-0.8cm] (qa) {QA Service\\(RAG + Gemini)};
\node[box, below=of gw] (meet) {Meetings\\Service};
\node[box, left=of meet] (auth) {Auth\\Service};
\node[box, below=2.6cm of asr] (db) {PostgreSQL\\+ pgvector};
\node[box, below=2.6cm of qa] (redis) {Redis\\(streams + tokens)};

\draw[arrow] (user) -- (fe);
\draw[arrow] (fe) -- (gw);
\draw[arrow] (gw) -- (asr);
\draw[arrow] (gw) -- (qa);
\draw[arrow] (gw) -- (meet);
\draw[arrow] (gw) -- (auth);
\draw[arrow] (asr) -- (db);
\draw[arrow] (asr) -- (redis);
\draw[arrow] (qa) -- (db);
\draw[arrow] (meet) -- (db);
\draw[arrow] (auth) -- (redis);
\end{tikzpicture}
\caption{High-level \meetai{} system context.}
\label{fig:high-level}
\end{figure}

\section{Thesis Organisation}
\label{sec:thesis-organisation}

The remainder of this thesis is organised as follows.

\begin{description}[itemsep=0.3em,leftmargin=4cm,style=nextline]
  \item[\Cref{ch:background}] reviews the relevant literature on ASR,
  speaker diarization, retrieval-augmented generation, and code-switched
  speech processing, and identifies the research gaps addressed by this
  thesis.
  \item[\Cref{ch:requirements}] presents the stakeholder analysis, the
  functional and non-functional requirements, and the data requirements
  of the proposed system.
  \item[\Cref{ch:architecture}] documents the system architecture and
  component design, including the high-level architecture, sequence
  diagrams for the key flows, the AI pipeline architecture, and the
  technology stack.
  \item[\Cref{ch:ml-methodology}] details the machine-learning
  methodology, including data preprocessing, the local-agreement
  streaming commit policy, the LoRA-based ASR fine-tuning strategy, the
  speaker-diarization integration, and the realtime RAG pipeline.
  \item[\Cref{ch:experiments}] describes the experimental setup,
  baselines, evaluation metrics, ablation studies, and statistical
  significance testing plan.
  \item[\Cref{ch:results}] presents and discusses the experimental
  results, including the model evaluation and the system evaluation, and
  discusses performance breakdown and failure modes.
  \item[\Cref{ch:conclusion}] concludes with a summary of contributions,
  critical reflection, and recommendations for future research.
  \item[\Cref{app:ethics,app:repro,app:genai}] provide the mandatory
  appendices on AI ethics, reproducibility, and generative-AI
  disclosure.
\end{description}

% =============================================================================
% Chapter 2: Literature Review (placeholder)
% =============================================================================

\chapter{Literature Review}
\label{ch:background}

\todo{Write Chapter 2: Literature Review per the guide's KSI requirements.}

\section{Domain Background}
\todo{Sub-sections 2.1 -- 2.4 will be filled in: domain background, SOTA software
solutions, survey of approaches (ASR, diarization, RAG), and identification of
research gaps.}

\section{State-of-the-Art in Software Solutions}
\todo{Survey existing meeting-assistant software (Otter, Fireflies, Granola,
Whisper-based local solutions) and their limitations for code-switched
Egyptian Arabic.}

\section{Survey of Approaches}
\todo{Three sub-sections: ASR architectures, speaker diarization, and
retrieval-augmented generation.}

\section{Identification of Research Gaps}
\todo{Summarize the gap addressed by this thesis: code-switched Egyptian
Arabic + realtime temporal-speaker-aware QA in a single integrated system.}

% =============================================================================
% Chapter 3: Analysis and Requirements Engineering (placeholder)
% =============================================================================

\chapter{Analysis and Requirements Engineering}
\label{ch:requirements}

\todo{Write Chapter 3 per the guide's KSI requirements.}

\section{Stakeholder Analysis}
\todo{Identify primary stakeholders (end users, team leads, project managers,
admins) and their pain points.}

\section{Functional Requirements}
\todo{Use case model. Include UML use case diagram and table mapping each
requirement to expected ML / system output.}

\section{Non-Functional Requirements}
\todo{Latency, scalability, robustness, security, privacy, accessibility.}

\section{Data Requirements}
\todo{Data sources, volumetric requirements, diversity requirements,
annotation standards. Define noise tolerance and class balance.}

% =============================================================================
% Chapter 4: System Architecture and Design (placeholder)
% =============================================================================

\chapter{System Architecture and Design}
\label{ch:architecture}

\todo{Write Chapter 4 per the guide's KSI requirements: high-level
architecture, component design, AI pipeline architecture, and technology
stack selection.}

\section{High-Level System Architecture}
\todo{System context diagram (already shown in Chapter 1, but expand here).}

\section{Component Design}
\todo{Class diagrams and sequence diagrams for the major flows (login,
start-meeting, streaming, QA, complete-meeting).}

\section{The AI Pipeline Architecture}
\todo{Sub-sections on data ingestion, preprocessing, model training,
inference, and API integration.}

\section{Technology Stack Selection}
\todo{Justify each technology: Node.js for gateway, Python for AI,
PostgreSQL+pgvector, Redis, Docker, Vite+React frontend.}

% =============================================================================
% Chapter 5: ML Methodology and Implementation (placeholder)
% =============================================================================

\chapter{ML Methodology and Implementation}
\label{ch:ml-methodology}

\todo{Write Chapter 5 per the guide's KSI requirements.}

\section{Data Preprocessing Techniques}
\todo{Audio normalization, voice activity detection, code-switch tokenization,
train/val/test split strategy.}

\section{Model Selection and Justification}
\todo{Whisper baseline choice, LoRA rationale, Sortformer (no fine-tune)
choice, sentence-transformer embedding choice, Gemini for topic/QA.}

\section{Feature Engineering and Selection}
\todo{Mel spectrogram features, text normalization, code-switch handling.}

\section{Hyperparameter Optimisation Strategy}
\todo{LoRA rank, learning rate, batch size, optimizer choice, scheduler. TBD
values will be filled in from the experiment log.}

\section{Software Implementation Details}
\todo{Algorithms, design patterns, and the local-agreement word-by-word
streaming commit policy. Code references to aiengine/audio/asr/online.py.}

% =============================================================================
% Chapter 6: Experimental Design and Evaluation (placeholder)
% =============================================================================

\chapter{Experimental Design and Evaluation}
\label{ch:experiments}

\todo{Write Chapter 6 per the guide's KSI requirements.}

\section{Experimental Setup}
\todo{Hardware (GPU/CPU/RAM), software (OS, Python, CUDA), versions,
deterministic seeds.}

\section{Baselines and Comparative Benchmarks}
\todo{At least two credible baselines (off-the-shelf Whisper, alternative
streaming ASR, alternative RAG).}

\section{Evaluation Metrics}
\todo{ASR: WER, CER, code-switch slice metrics. Diarization: DER, speaker
attribution accuracy. Retrieval: Recall@k, MRR, nDCG. QA: faithfulness and
groundedness. System: latency, throughput.}

\section{Ablation Studies}
\todo{With/without LoRA, with/without local-agreement, with/without
diarization, with/without topic indexing.}

\section{Statistical Significance Testing}
\todo{Test choice (paired bootstrap, McNemar, etc.), confidence intervals,
significance threshold.}

\section{Model Evaluation}
\todo{Fill the TBD tables in \cref{sec:model-eval} (Chapter 7) with values
from the experiment log.}

\section{System Evaluation}
\todo{Fill the TBD tables in \cref{sec:system-eval} (Chapter 7) with values
from the experiment log.}

% =============================================================================
% Chapter 7: Results and Discussion (placeholder)
% =============================================================================

\chapter{Results and Discussion}
\label{ch:results}

\todo{Write Chapter 7 per the guide's KSI requirements.}

\section{Analysis of Experimental Results}
\todo{Discuss the model evaluation results from \cref{sec:model-eval}.}

\section{Performance Breakdown}
\todo{Edge cases, code-switch slices, minority classes, overlapping
speakers, ambiguous time references.}

\section{Implications for Software Engineering Practice}
\todo{Connect ML performance back to the original SE problem; discuss
trade-offs between model complexity and computational cost.}

\section{Limitations of the Current Approach}
\todo{Honest discussion of limitations (e.g., small fine-tune dataset,
single-dialect, no overlap handling, English-only QA generation).}

% ---------------------------------------------------------------------------
% 7.5 Model Evaluation (TBD)
% ---------------------------------------------------------------------------
\section{Model Evaluation}
\label{sec:model-eval}

\subsection{ASR Evaluation}
\todo{Fill the metrics below from the experiment log.}

\begin{table}[H]
\centering
\caption{ASR quality on the code-switched Egyptian Arabic test set. TBD
values will be replaced with measured metrics. Best per row in
\textbf{bold}. BS = baseline (off-the-shelf Whisper), LoRA = fine-tuned
model, $k$ = local-agreement threshold.}
\label{tab:asr-metrics}
\begin{tabularx}{\textwidth}{l X X X X}
\toprule
\textbf{Model} & \textbf{Setting} & \textbf{WER} & \textbf{CER} & \textbf{Notes} \\
\midrule
BS, $k=1$     & \tbd & \tbd & \tbd & streaming, no commit \\
BS, $k=3$     & \tbd & \tbd & \tbd & local-agreement \\
LoRA, $k=1$   & \tbd & \tbd & \tbd & \tbd \\
LoRA, $k=3$   & \tbd & \tbd & \tbd & \tbd \\
LoRA, $k=5$   & \tbd & \tbd & \tbd & \tbd \\
\bottomrule
\end{tabularx}
\end{table}

\begin{table}[H]
\centering
\caption{ASR quality broken down by language slice. TBD.}
\label{tab:asr-by-slice}
\begin{tabularx}{\textwidth}{l X X X}
\toprule
\textbf{Model} & \textbf{Egyptian Arabic WER} & \textbf{English WER} & \textbf{Code-Switch WER} \\
\midrule
BS             & \tbd & \tbd & \tbd \\
LoRA (ours)    & \tbd & \tbd & \tbd \\
\bottomrule
\end{tabularx}
\end{table}

\subsection{Diarization Evaluation}

\begin{table}[H]
\centering
\caption{Speaker-diarization performance. TBD.}
\label{tab:diarization-metrics}
\begin{tabularx}{\textwidth}{l X X X}
\toprule
\textbf{Setting} & \textbf{DER} & \textbf{Speaker Attribution Acc.} & \textbf{Notes} \\
\midrule
Pretrained Sortformer       & \tbd & \tbd & \tbd \\
With sortformer + alignment & \tbd & \tbd & \tbd \\
\bottomrule
\end{tabularx}
\end{table}

\subsection{Retrieval and QA Evaluation}

\begin{table}[H]
\centering
\caption{Retrieval and QA quality. TBD.}
\label{tab:retrieval-qa-metrics}
\begin{tabularx}{\textwidth}{l X X X X}
\toprule
\textbf{Setting} & \textbf{Recall@5} & \textbf{MRR} & \textbf{nDCG@10} & \textbf{Answer Faithfulness} \\
\midrule
Chunk-only              & \tbd & \tbd & \tbd & \tbd \\
Topic + chunk fusion    & \tbd & \tbd & \tbd & \tbd \\
With speaker filter     & \tbd & \tbd & \tbd & \tbd \\
With time filter        & \tbd & \tbd & \tbd & \tbd \\
\bottomrule
\end{tabularx}
\end{table}

% ---------------------------------------------------------------------------
% 7.6 System Evaluation (TBD)
% ---------------------------------------------------------------------------
\section{System Evaluation}
\label{sec:system-eval}

\subsection{Streaming ASR Latency}

\begin{table}[H]
\centering
\caption{Streaming ASR latency and real-time factor. TBD.}
\label{tab:asr-latency}
\begin{tabularx}{\textwidth}{l X X X}
\toprule
\textbf{Setting} & \textbf{First-Token Latency} & \textbf{Word Commit Latency} & \textbf{RTF} \\
\midrule
$k=1$ (no local agreement)  & \tbd & \tbd & \tbd \\
$k=3$ (local agreement)     & \tbd & \tbd & \tbd \\
$k=5$ (strict)              & \tbd & \tbd & \tbd \\
\bottomrule
\end{tabularx}
\end{table}

\subsection{End-to-End QA Latency}

\begin{table}[H]
\centering
\caption{End-to-end QA latency for temporal-speaker queries. TBD.}
\label{tab:qa-latency}
\begin{tabularx}{\textwidth}{l X X}
\toprule
\textbf{Query Type} & \textbf{P50 Latency} & \textbf{P95 Latency} \\
\midrule
Speaker-only                 & \tbd & \tbd \\
Time-only                    & \tbd & \tbd \\
Speaker + time               & \tbd & \tbd \\
Speaker + time + paraphrase  & \tbd & \tbd \\
\bottomrule
\end{tabularx}
\end{table}

\subsection{Throughput, Concurrency, and Resource Profile}

\begin{table}[H]
\centering
\caption{Throughput, concurrency behavior, and resource profile. TBD.}
\label{tab:throughput}
\begin{tabularx}{\textwidth}{l X X X X}
\toprule
\textbf{Concurrent Sessions} & \textbf{Mean CPU \%} & \textbf{GPU Mem (GB)} & \textbf{RAM (GB)} & \textbf{Drop Rate} \\
\midrule
1     & \tbd & \tbd & \tbd & \tbd \\
4     & \tbd & \tbd & \tbd & \tbd \\
8     & \tbd & \tbd & \tbd & \tbd \\
16    & \tbd & \tbd & \tbd & \tbd \\
\bottomrule
\end{tabularx}
\end{table}

\subsection{Stability over Long Sessions}

\todo{Report transcript-regression rate, speaker-label-swap rate, and
session-disconnect rate for meetings of 30 / 60 / 90 minutes. TBD.}

% =============================================================================
% Chapter 8: Conclusion and Future Work (placeholder)
% =============================================================================

\chapter{Conclusion and Future Work}
\label{ch:conclusion}

\todo{Write Chapter 8 per the guide's KSI requirements.}

\section{Summary of Contributions}
\todo{Bullet list of contributions: corpus, LoRA-fine-tuned ASR,
local-agreement commit policy, RAG pipeline, temporal-speaker QA, full
integrated system.}

\section{Critical Reflection}
\todo{What worked, what did not work, what we would do differently.}

\section{Recommendations for Future Research}
\todo{Overlap handling, broader dialect coverage, on-device deployment,
explicit evaluation of hallucination in QA answers, integration with
calendar actions, etc.}


% =============================================================================
% APPENDICES
% =============================================================================
\appendix
% =============================================================================
% Appendix A: AI Ethics & Bias Statement
% =============================================================================

\chapter{AI Ethics \& Bias Statement}
\label{app:ethics}

\todo{Write Appendix A per the guide's KSI requirements.}

\section{Data Fairness}
\todo{Describe demographic analysis of the curated corpus: gender, age
range, regional dialect variation, accent diversity. Identify potential
biases.}

\section{Mitigation}
\todo{Steps taken to balance the dataset (e.g., stratified sampling, target
dialect representation, targeted data collection for underrepresented
groups).}

\section{Transparency}
\todo{Disclose the model's black-box nature and any explainability
techniques used (e.g., word-level confidence, retrieval-evidence
grounding). Discuss dual-use and misuse risks and the guardrails in
place.}

\section{Ethical Risk Assessment}
\todo{Identify potential dual-use / misuse cases (e.g., unauthorized
recording, surveillance) and the technical and policy mitigations.}

% =============================================================================
% Appendix B: Reproducibility Report
% =============================================================================

\chapter{Reproducibility Report}
\label{app:repro}

\todo{Write Appendix B per the guide's KSI requirements: environment,
containerization, deterministic seeds, hardware specs, hyperparameter log.}

\section{Environment}
\todo{Link to / include \texttt{requirements.txt} (Python) and
\texttt{package.json} (Node), plus the root \texttt{ai/.env.example} and
\texttt{server/.env.example}.}

\section{Containerization}
\todo{Describe the \texttt{docker-compose} setup and how a reviewer can
recreate the environment with a single command.}

\section{Deterministic Seeds}
\todo{List all random seeds used for data splitting, weight initialization,
embedding model inference, and any stochastic decoding.}

\section{Hardware Specs}
\todo{GPU model and memory, CPU model and core count, RAM, disk,
network. Include the configuration used for both fine-tuning and
inference benchmarks.}

\section{Hyperparameter Log}
\todo{Table of all tuned parameters (not only the winning ones) and the
search space. TBD: fill from the experiment log.}

% =============================================================================
% Appendix C: Generative AI (GenAI) Disclosure
% =============================================================================

\chapter{Generative AI (GenAI) Disclosure}
\label{app:genai}

\todo{Write Appendix C per the guide's KSI requirements.}

\section{Tool Inventory}
\todo{List all LLMs and AI assistants used (e.g., Claude, Gemini,
GitHub Copilot).}

\section{Usage Scope}
\todo{State explicitly the scope of AI usage: code generation, debugging,
drafting prose, diagram generation, and any other categories.}

\section{Verification}
\todo{Confirm that all AI-generated content was audited, corrected, and
verified by the project team before inclusion.}

\section{Critical Audit Log}
\todo{Include examples of where an LLM provided an incorrect or
hallucinated answer and how the project team corrected it.}


% =============================================================================
% REFERENCES / BIBLIOGRAPHY
% =============================================================================
\backmatter
\bibliography{references}

\end{document}
